% 高中政治：辩证法入门与进阶 笔记
% 使用 mathnote-preamble.tex 模板，建议用 XeLaTeX 编译两次

\documentclass[a4paper,11pt]{ctexart}
\input{../mathnote-preamble.tex}

% 元信息配置（用于页眉、书签与封面）
\renewcommand{\notetitle}{高中政治：辩证法入门与进阶}
\renewcommand{\notesubtitle}{联系·矛盾·发展观的系统学习}
\renewcommand{\noteauthor}{自学笔记示例}
\renewcommand{\notedate}{\today}
\renewcommand{\noteversion}{v1.0}

\begin{document}

% 封面
\thispagestyle{empty}
\PageTag[secondary]{高中政治·辩证法}
\setcounter{page}{1}

\MathNoteEnableReviewStamp
\MathNoteEnablePrint

\noindent
\begin{minipage}[t][0.7\textheight][c]{\textwidth}
  \raggedleft
  \vspace*{\fill}
  {\sffamily\bfseries\Huge\color{accent}\notetitle\par}
  \vspace{0.6em}
  {\Large\color{inkgray}\notesubtitle\par}
  \vspace{0.5em}
  {\large\color{inkgray}\noteauthor\par}
  \vspace{0.4em}
  {\normalsize\color{inkgray}\noteversion\quad|\quad\notedate\par}
  \vspace*{\fill}
\end{minipage}

\begin{center}
\begin{minipage}{0.84\textwidth}
  \textcolor{inkgray}{\sffamily\small 学习导航}\\[0.4em]
  \begin{focuspoints}
    \item 从\keyword{联系与发展}的整体视角出发，理解辩证法区别于日常“想当然”的思维方式；
    \item 掌握\keyword{矛盾分析法}、\keyword{量变—质变规律}、\keyword{否定之否定规律}等核心工具；
    \item 通过简明\keyword{框架图}和\keyword{对比表}理清易混概念，形成可复用的答题模板；
    \item 结合\keyword{学习生活实例}与\keyword{高考题型}，学会把抽象哲学语言转化为具体分析步骤；
    \item 搭建“\keyword{读题—找点—组织语言}”的一条龙流程，实现自学与自我纠错。
  \end{focuspoints}
\end{minipage}
\end{center}

\clearpage
\thispagestyle{plain}
\phantomsection
\tableofcontents
\newpage

%==========================
\section{导学与整体结构}
%==========================

\begin{summarybox}{为什么要学辩证法？}
在高中政治中，辩证法既是\keyword{世界观}（怎么看世界），也是\keyword{方法论}（怎么分析问题）。理解辩证法，可以帮助我们：
\begin{itemize}
  \item 看见事物之间的\keyword{联系}，避免“只盯一面”的片面性；
  \item 把握事物\keyword{发展变化}的方向与条件，而不只看静止结果；
  \item 通过\keyword{矛盾分析}找到关键点，从而更有针对性地解决问题；
  \item 在高考题中用规范语言回答“\inlinehint{为什么}”“\inlinehint{怎么办}”类问题。
\end{itemize}
\end{summarybox}

\begin{roadmap}
  \RoadmapStep{建立辩证法的整体印象}
  \RoadmapStep{掌握联系与发展的基本观点}
  \RoadmapStep{重点突破矛盾分析法}
  \RoadmapStep{理解量变—质变与否定之否定}
  \RoadmapStep{构建答题模板与应用步骤}
\end{roadmap}

\inlinehint{建议学习顺序：先把每节的“定义盒 + 对比表 + 实例”吃透，再尝试用本章的答题模板分析课本或真题。}

%==========================
\section{辩证法的基本视角}
%==========================

\subsection{什么是辩证法？}

\begin{definitionbox}{辩证法的基本含义}
\keyword{辩证法}是一种关于世界\keyword{普遍联系}和\keyword{永恒发展}的学说，也是人们认识和改造世界的\keyword{科学方法}。用简明话概括：
\begin{itemize}
  \item 世界不是静止的，而是在\keyword{运动变化}中发展；
  \item 事物不是孤立的，而是处在\keyword{普遍联系}之中；
  \item 发展不是一帆风顺的，而是通过\keyword{矛盾运动}实现的。
\end{itemize}
\end{definitionbox}

\begin{examplebox}{从“备考”看两种思维方式}
\textbf{情境：}同样面对高考复习，两位同学的想法不同：
\begin{itemize}
  \item 甲：只盯着分数和排名，认为“我现在基础差，肯定追不上”；
  \item 乙：既看到当前差距，又分析\inlinehint{时间、方法、情绪、资源}等多个因素，制定分阶段提升计划。
\end{itemize}
\textbf{分析：}
\begin{itemize}
  \item 甲忽视了学习过程中的\keyword{发展性}和多种\keyword{联系因素}，思维更接近\inlinehint{静止、孤立}的看法；
  \item 乙从多种条件出发，把当前水平看作发展过程中的\keyword{一个阶段}，体现了\keyword{联系和发展的辩证观点}。
\end{itemize}
\end{examplebox}

\subsection{辩证法与形而上学的对比}

\begin{summarybox}{两种世界观的对比框架}
在高中政治中，常用\keyword{辩证法}与\keyword{形而上学}进行对比。可以记成四组关键词：
\begin{focuspoints}
  \item \keyword{整体 vs. 孤立}：看不看整体联系；
  \item \keyword{发展 vs. 静止}：承不承认变化过程；
  \item \keyword{矛盾 vs. 单一}：看到差异和冲突还是只看统一；
  \item \keyword{转化 vs. 固化}：承不承认条件改变可以带来方向改变。
\end{focuspoints}
\end{summarybox}

\begin{center}
\begin{tabularx}{0.96\textwidth}{lXX}
\toprule
视角 & 辩证法思维 & 形而上学思维 \\
\midrule
对事物状态的看法 & 强调事物\keyword{处在变化发展之中}，任何状态都是过程中的一个阶段 & 把事物视为\inlinehint{固定不变}，把某一时刻当成“永远如此” \\
对事物关系的看法 & 强调\keyword{普遍联系}，从整体中看各部分作用与影响 & 只盯住某一局部或某一个因素，忽视其与环境的联系 \\
对矛盾的态度 & 承认并重视\keyword{矛盾}，善于找到\keyword{主要矛盾和矛盾的主要方面} & 不是否认矛盾，就是只看到矛盾的一方，做“非黑即白”的判断 \\
方法论要求 & 要求\keyword{具体问题具体分析}，在联系和发展中寻找解决方案 & 喜欢套用固定结论，容易陷入\inlinehint{绝对化、一刀切、简单类比} \\
\bottomrule
\end{tabularx}
\end{center}

\begin{center}
\begin{tikzpicture}[mathnote lines, node distance=2.6cm]
  \node[draw,rounded corners=3pt,fill=surface,inner sep=6pt] (all) {\sffamily\small 世界是\keyword{普遍联系}和\keyword{永恒发展}的整体};
  \node[draw,rounded corners=3pt,fill=surface,below left of=all,xshift=-0.6cm] (dial) {\sffamily\small 辩证法：\\联系 + 发展 + 矛盾};
  \node[draw,rounded corners=3pt,fill=surface,below right of=all,xshift=0.6cm] (meta) {\sffamily\small 形而上学：\\孤立 + 静止 + 片面};
  \draw[->,accent,thick] (all.south west) -- (dial.north);
  \draw[->,secondary,thick] (all.south east) -- (meta.north);
\end{tikzpicture}
\captionof{figure}{两种基本世界观的示意图}
\end{center}

%==========================
\section{联系与发展的观点}
%==========================

\subsection{联系的观点}

\begin{definitionbox}{联系的含义与特点}
\keyword{联系}是指事物内部各要素之间以及事物之间、事物与环境之间的\keyword{相互影响和相互制约}。其主要特点：
\begin{itemize}
  \item \keyword{普遍性}：世界上一切事物都处在一定联系之中；
  \item \keyword{多样性}：联系有直接与间接、内部与外部、必然与偶然等不同形式；
  \item \keyword{条件性}：不同条件下，联系的\inlinehint{强弱、方向和方式}可能发生变化。
\end{itemize}
\end{definitionbox}

\begin{center}
\begin{tabularx}{0.96\textwidth}{lXX}
\toprule
联系类型 & 简要说明 & 示例 \\
\midrule
内部联系 & 同一事物内部诸要素之间的联系 & 一次考试成绩由\keyword{基础、临场状态、审题习惯}等多个内部因素共同决定 \\
外部联系 & 事物与外部环境之间的联系 & 学习效率还受\keyword{家庭氛围、学校资源、同伴影响}等外在条件制约 \\
必然联系 & 在一定条件下必然出现的、稳定的联系 & 长期坚持高质量练习\inlinehint{必然}会提高熟练度和正确率 \\
偶然联系 & 在偶然条件下形成、具有一定偶然性的联系 & 某次蒙对选择题只是\inlinehint{偶然}，不能说明整体水平提高 \\
\bottomrule
\end{tabularx}
\end{center}

\begin{examplebox}{分析“网课学习效果”的多重联系}
网课学习效果和哪些因素有关？尝试从不同类型的\keyword{联系}出发分析：
\begin{itemize}
  \item 内部联系：自身\keyword{自律程度、学习习惯、专注时间}等；
  \item 外部联系：网络环境、家人干扰、老师互动方式等；
  \item 必然联系：长期高质量的\keyword{听课 + 训练 + 复盘}会带来稳步提升；
  \item 偶然联系：一两节课\keyword{状态极好或极差}，不能简单外推为“我就这样了”。
\end{itemize}
此类题目作答时，可按“\keyword{指出联系} $\rightarrow$ \keyword{举例说明} $\rightarrow$ \keyword{提出建议}”的顺序组织语言。
\end{examplebox}

\subsection{发展的观点}

\begin{summarybox}{发展的含义与基本特征}
从辩证法角度看，\keyword{发展}是事物由低级到高级、由简单到复杂、由旧质到新质的\keyword{前进性变化}。其特点可以记为“三个不”：
\begin{itemize}
  \item \keyword{不是一帆风顺的}：发展包含曲折，前进中有反复；
  \item \keyword{不是自动完成的}：需要一定的\inlinehint{条件和主观努力}；
  \item \keyword{不是无方向的}：总体上具有\keyword{上升、前进}的趋势。
\end{itemize}
\end{summarybox}

\begin{center}
\begin{tikzpicture}[mathnote lines,xscale=1.2,yscale=0.8]
  \draw[->,thick] (0,0) -- (6,0) node[right] {时间};
  \draw[->,thick] (0,0) -- (0,4) node[above] {水平};
  \draw[accent,thick,smooth] plot coordinates {(0.3,0.6) (1,1.2) (1.8,1.0) (2.5,1.8) (3.3,2.1) (4,2.5) (4.7,2.3) (5.5,3.2)};
  \node[anchor=west,inkgray] at (3.1,3.0) {\scriptsize 曲折中前进的学习曲线};
\end{tikzpicture}
\captionof{figure}{“曲折上升”的发展示意图}
\end{center}

%==========================
\section{矛盾分析：从现象到本质}
%==========================

\subsection{矛盾及其基本规律}

\begin{definitionbox}{矛盾及其两种属性}
\keyword{矛盾}是指\keyword{事物内部或事物之间}既对立又统一的关系。矛盾具有：
\begin{itemize}
  \item \keyword{普遍性}：矛盾存在于一切事物中，贯穿事物发展全过程；
  \item \keyword{特殊性}：不同事物的矛盾各不相同，同一事物在不同发展阶段的矛盾也不同。
\end{itemize}
把握“\keyword{普遍性}—\keyword{特殊性}”关系，是实现\keyword{具体问题具体分析}的理论依据。
\end{definitionbox}

\begin{summarybox}{主要矛盾与主要方面}
\begin{itemize}
  \item \keyword{主要矛盾}：在众多矛盾中，对事物发展起\keyword{主导、决定}作用的矛盾；
  \item \keyword{次要矛盾}：处于从属地位、影响较小的其他矛盾；
  \item \keyword{矛盾的主要方面}：在每一对矛盾双方中处于\keyword{支配地位}的一方；
  \item \keyword{矛盾的次要方面}：处于\inlinehint{相对被动}的一方。
\end{itemize}
在分析复杂问题时，应学会\keyword{抓主要矛盾}，并在每一矛盾中注意\keyword{把握主要方面}。
\end{summarybox}

\begin{center}
\begin{tabularx}{0.96\textwidth}{lXX}
\toprule
矛盾层次 & 典型问题 & 分析提示 \\
\midrule
主要矛盾—次要矛盾 & 学习成绩下降，到底是\keyword{心态问题}还是\keyword{方法问题}为主？ & 先列出所有矛盾，再判断哪一矛盾\inlinehint{影响最大、作用最长远} \\
主要方面—次要方面 & 在“时间不够”和“效率不高”的矛盾中，哪一方更关键？ & 分析哪一方面的改变量更可能\keyword{改变整体状况} \\
\bottomrule
\end{tabularx}
\end{center}

\subsection{矛盾分析四步法}

\begin{summarybox}{矛盾分析的操作流程}
可以把“具体问题具体分析”固化为四步：
\begin{enumerate}
  \item \keyword{找矛盾}：问题中\inlinehint{对立的两方面}分别是什么？
  \item \keyword{划层次}：哪一对矛盾是\keyword{主要矛盾}？哪一方面是\keyword{主要方面}？
  \item \keyword{看条件}：在什么\inlinehint{具体条件}下，它们的地位和作用会发生变化？
  \item \keyword{提对策}：针对主要矛盾和主要方面提出\keyword{有侧重点}的解决方案。
\end{enumerate}
\end{summarybox}

\begin{center}
\begin{tikzpicture}[mathnote lines,node distance=1.9cm]
  \node[draw,rounded corners=3pt,fill=surface,inner sep=5pt] (step1) {\scriptsize 1. 找矛盾};
  \node[draw,rounded corners=3pt,fill=surface,right=of step1] (step2) {\scriptsize 2. 划层次};
  \node[draw,rounded corners=3pt,fill=surface,right=of step2] (step3) {\scriptsize 3. 看条件};
  \node[draw,rounded corners=3pt,fill=surface,right=of step3] (step4) {\scriptsize 4. 提对策};
  \draw[->,accent,thick] (step1) -- (step2);
  \draw[->,accent,thick] (step2) -- (step3);
  \draw[->,accent,thick] (step3) -- (step4);
\end{tikzpicture}
\captionof{figure}{矛盾分析四步法流程图}
\end{center}

\begin{examplebox}{用矛盾分析“玩手机与学习”的冲突}
\textbf{情境：}备考期间，你发现自己经常“想学习又忍不住刷手机”。尝试用矛盾分析四步法梳理：
\begin{enumerate}
  \item 找矛盾：\keyword{提高成绩的需要}与\keyword{即时娱乐的欲望}之间的矛盾；
  \item 划层次：当前阶段，高考时间临近，“提高成绩”是\keyword{主要矛盾的主要方面}；
  \item 看条件：无计划、无节制的娱乐会严重\inlinehint{挤占学习时间}，但适度休息有利于恢复精力；
  \item 提对策：设置固定“刷手机时段”，并用\inlinehint{番茄钟、应用锁}等外部工具帮助控制，将主要精力集中到学习上。
\end{enumerate}
答题时，可以用“首先—其次—再次—最后”或“原因—影响—对策”的结构，把上述分析用规范语言表述出来。
\end{examplebox}

%==========================
\section{量变—质变与否定之否定}
%==========================

\subsection{量变与质变}

\begin{definitionbox}{量变与质变的基本概念}
\begin{itemize}
  \item \keyword{量变}：在事物的\inlinehint{数量、程度、范围}等方面发生的\keyword{渐进性变化}，一般比较缓慢、连续；
  \item \keyword{质变}：当量变发展到一定程度时，引起事物\keyword{性质、结构、状态}的\keyword{根本变化}。
\end{itemize}
量变和质变是\keyword{互相渗透、互相转化}的：没有量变的积累就不会发生质变，质变又为新的量变开辟道路。
\end{definitionbox}

\begin{center}
\begin{tabularx}{0.96\textwidth}{lXX}
\toprule
类型 & 关键词 & 示例 \\
\midrule
量变 & 渐进、积累、数量、程度 & 每天多做几道错题整理，多记几个答题句式 \\
质变 & 突破、飞跃、性质改变 & 从“看不懂���”到“能独立完整作答”，成绩段位明显上升 \\
\bottomrule
\end{tabularx}
\end{center}

\begin{examplebox}{从“刷题”到“真正掌握”的量变—质变}
刚开始刷题时，你可能\inlinehint{错误很多、速度很慢}，但如果坚持：
\begin{itemize}
  \item 每道错题都\keyword{纠正原因}，总结对应知识点；
  \item 定期回顾错题本，观察错误类型是否在\inlinehint{减少或变化}；
  \item 按章节和题型整理\keyword{解题步骤和常用句式}。
\end{itemize}
这些看似不起眼的\keyword{量变}积累到一定程度，就会表现为：
\begin{itemize}
  \item 看到题目后能迅速\keyword{判断考点和设问类型}；
  \item 作答用语更加\keyword{规范、全面、有针对性}；
  \item 分数出现\keyword{阶段性跃升}——这就是学习中的\keyword{质变}。
\end{itemize}
\end{examplebox}

\subsection{否定之否定}

\begin{summarybox}{否定之否定规律的简要理解}
\keyword{否定之否定}不是简单的“推翻”，而是\keyword{既克服又保留}的扬弃过程。可以记成三个关节点：
\begin{enumerate}
  \item \keyword{第一次否定}：旧事物走向灭亡，新事物产生；
  \item \keyword{第二次否定}：新事物发展中的更高级阶段，对原来的新事物进行\inlinehint{改造和提升}；
  \item \keyword{结果形态}：既不同于原来的旧事物，又保留了其中的\keyword{合理因素}，形成“\inlinehint{螺旋式上升、波浪式前进}”的过程。
\end{enumerate}
\end{summarybox}

\begin{examplebox}{从“旧笔记”到“新框架”的否定之否定}
\begin{enumerate}
  \item 第一次否定：你发现原来的课堂笔记\inlinehint{杂乱、难以复习}，于是决定重新整理辩证法内容，建立更清晰的结构；
  \item 第二次否定：在使用新笔记一段时间后，你根据\keyword{错题反馈和模拟考试}进一步调整结构，增补“答题模板”和“高频错误”板块；
  \item 结果形态：最终形成一套\keyword{结构清晰、针对性强、可循环使用}的专用笔记，本身就是对过去两版笔记的\keyword{扬弃和提升}。
\end{enumerate}
\end{examplebox}

%==========================
\section{辩证法在高考答题中的应用}
%==========================

\subsection{常见设问类型与对应工具}

\begin{summarybox}{“看题型选工具”的思路}
面对辩证法相关题目，可以先判断设问属于哪一种，再选用对应的\keyword{工具与句式}：
\end{summarybox}

\begin{center}
\begin{tabularx}{0.96\textwidth}{lXX}
\toprule
设问类型 & 对应辩证法工具 & 作答思路提示 \\
\midrule
“如何全面认识……？” & 联系观点、两点论与重点论 & 从\keyword{多角度列点}，再突出\keyword{主要方面或关键因素} \\
“如何看待利与弊？” & 矛盾分析、对立统一 & 从\keyword{正反两方面}分析，强调“寓利于弊、趋利避害” \\
“如何正确处理……关系？” & 主要矛盾与次要矛盾 & 说明\keyword{谁是主要矛盾}，强调“抓住主要矛盾同时兼顾其他矛盾” \\
“怎样做到稳中求进？” & 量变—质变规律 & 先讲\keyword{打好基础、循序渐进}，再讲\keyword{适时把握质变时机} \\
“如何持续优化……？” & 否定之否定规律 & 强调在\keyword{阶段性总结中扬弃不足}，实现\keyword{螺旋式上升} \\
\bottomrule
\end{tabularx}
\end{center}

\subsection{示例：分析一段材料}

\begin{examplebox}{用辩证思维分析学习计划}
\textbf{材料：}某班组织“百日学习打卡”活动，有同学认为“只要每天坚持刷题就一定能提高成绩”；也有人认为“活动太累，影响正常生活，没有必要参加”。请运用辩证法知识加以分析。

\textbf{示例思路：}
\begin{enumerate}
  \item 联系与发展：既要看到“刷题活动”对\keyword{熟练度、氛围、自律}等积极影响，也要看到\inlinehint{身体健康、心理压力、活动质量}等因素；
  \item 矛盾分析：当前主要矛盾是“\keyword{提高学习质量}与\keyword{避免形式主义}”之间的矛盾，应把握“\keyword{提高质量}”这一主要方面；
  \item 量变—质变：合理的每日积累有利于形成\keyword{质的飞跃}，但一味追求数量可能\inlinehint{降低效率、积累负面量变}；
  \item 否定之否定：可以通过\keyword{调整活动内容和节奏}，保留打卡的\inlinehint{激励和监督}作用，否定其中\inlinehint{机械刷题、疲劳战}等不合理部分，使活动不断优化。
\end{enumerate}

\textbf{规范作答示例要点：}
\begin{itemize}
  \item 观点句：坚持用\keyword{联系的、发展的、全面的}观点看问题；
  \item 分析句：既要\keyword{肯定}有利于提高学习质量的积极作用，又要\keyword{看到}可能带来的负面影响；
  \item 对策句：在\keyword{抓住提高质量这一主要方面}的前提下，\keyword{优化活动设计}，实现\keyword{量变积累到质变飞跃}的良性发展。
\end{itemize}
\end{examplebox}

%==========================
\section{本章小结与自我检测}
%==========================

\begin{summarybox}{核心知识网络回顾}
可以按下面的\keyword{“三层网络”}自查是否掌握：
\begin{itemize}
  \item \textbf{观念层（知道什么）}：能否简明说出\keyword{联系、发展、矛盾、量变、质变、否定之否定}等基本概念及其关系；
  \item \textbf{工具层（会用什么）}：遇到设问时，是否能迅速判断用\keyword{联系观点、矛盾分析、量变—质变}还是\keyword{否定之否定}来作答；
  \item \textbf{应用层（能做到什么）}：是否能用\keyword{规范政治语言}分析自身学习、班级活动或社会热点。
\end{itemize}
\end{summarybox}

\begin{focuspoints}
  \item 尝试用“矛盾分析四步法”写一段关于\keyword{手机使用}或\keyword{熬夜学习}的分析；
  \item 任选一道最近做过的高考政治题，标注其中涉及的\keyword{辩证法知识点}并改写自己的答案；
  \item 定期\keyword{更新本笔记}，把新的例子和答题句式补充进来，使之成为你个人的\keyword{辩证思维工具箱}。
\end{focuspoints}

\end{document}

